\bookchapter{Chapter 30: Screen Mirroring, Presentations, and Car Connectivity}

\chapterintro{Your phone screen is at most 7 inches, but your TV is 55 inches and your projector is 100 inches. Casting your phone screen to a larger display is a completely different experience. This chapter covers three high-frequency scenarios: screen casting, presentations, and car connectivity — helping you extend your phone's capabilities to bigger screens.}

\booksection{I. Three Wireless Casting Camps — Don't Mix Them Up}

Wireless casting all looks like “phone screen to TV,” but there are three underlying technologies that are mutually incompatible:

\begin{center}
\begin{tabularx}{\linewidth}{|X|X|X|}
\hline
\textbf{Technology} & \textbf{Who Uses It} & \textbf{Protocol} \tabularnewline
\hline
AirPlay & iPhone/iPad/Mac & Apple ecosystem, smoothest experience. Requires Apple TV or AirPlay-enabled TV \tabularnewline
Miracast & Android phones & Wi-Fi Direct connection, no router needed. Android phones' "wireless casting" uses this by default \tabularnewline
DLNA & All platforms & Only pushes media files (video/music/photos), doesn't mirror screen. Good for pushing video to TV for playback \tabularnewline
\hline
\end{tabularx}
\end{center}

\booksubsection{Operation methods for each platform:}

\begin{itemize}
  \item \textbf{iPhone}: Control Center > Screen Mirroring > Select your TV/Apple TV. If the TV doesn't support AirPlay, you can buy an Apple TV box or an AirPlay-enabled TV stick.
  \item \textbf{Android}: Swipe down the notification panel > Find "Wireless Cast" / "Smart Cast" / "Multi-Screen Interaction" button > Select TV. Phone and TV must be on the same Wi-Fi.
  \item \textbf{Universal Solution}: If the TV is too old to support any wireless casting, plug a Chromecast or Xiaomi Box / Huawei Box into the TV's HDMI port to receive casts.
\end{itemize}

\booksection{II. Wired Casting — More Stable, Zero Latency}

Wireless casting is convenient but has latency, making gaming and presentations unpleasant. Wired casting solves this:

\begin{itemize}
  \item \textbf{iPhone}: Requires a Lightning/USB-C to HDMI adapter cable (about 100-200 yuan). Plug in, and the phone screen outputs directly to the TV/projector's HDMI port.
  \item \textbf{Android}: Some phones support USB-C to HDMI direct connection (requires DP Alt Mode support; most flagships support it). Just buy a USB-C to HDMI cable.
  \item \textbf{Meeting Room Scenario}: If you frequently use your phone for presentations, keep an HDMI adapter cable in your bag. Meeting room TVs/projectors almost all have HDMI ports — just plug in and go, far more reliable than fiddling with wireless casting.
\end{itemize}

\booksection{III. Phone Becomes Computer — Desktop Mode}

Some Android phones support “Desktop Mode” — when connected to a monitor, the phone outputs a computer-like desktop interface that can be operated with mouse and keyboard:

\begin{itemize}
  \item \textbf{Samsung DeX}: Samsung flagship's signature feature. Connecting a monitor brings up a desktop interface with multi-window support, the closest experience to a real computer. With Bluetooth keyboard and mouse, you can travel without a laptop.
  \item \textbf{Huawei Desktop Mode}: Huawei Mate series / Pura series also has similar functionality when connected to a monitor. Supports multi-window, file management, and office apps.
  \item \textbf{Motorola Ready For}: Moto flagships also support it, with features similar to DeX.
  \item \textbf{What is RAW}: Light office work (email, document editing), PPT presentations, watching videos on big screen. Not suitable for heavy Excel work or professional design.
\end{itemize}

\booksection{IV. Car Connectivity — The Right Way to Use Your Phone While Driving}

Operating your phone while driving is illegal, but navigation, music, and calls are essential. Car connectivity systems let you use your phone's core features on the car's infotainment screen, so the phone can be set aside:

\begin{itemize}
  \item \textbf{Apple CarPlay}: Top choice for iPhone users. Connect to the car system via USB or wirelessly, and the car screen shows a streamlined iPhone interface — navigation (Amap/Baidu Maps), music, phone, messages, podcasts. Siri voice control is entirely hands-free.
  \item \textbf{Android Auto}: Android users' counterpart. Features similar to CarPlay, supporting Google Maps, Spotify, etc. Some features are unavailable in China due to Google service restrictions.
  \item \textbf{Huawei HiCar}: Huawei's proprietary solution, widely supported by domestic car models. Beyond navigation and music, it also supports Huawei video calls, parking camera fusion, and other unique features.
  \item \textbf{Baidu CarLife}: Cross-platform solution, works with both iOS and Android. Fewer features but broad compatibility — many older car models only support CarLife.
\end{itemize}

\begin{tcbwarn}{Security Reminder}
Regardless of which car connectivity solution you use, always rely on voice control while driving. Set your navigation destination and music playlist before starting to drive — never operate the car screen while in motion.

\end{tcbwarn}

\booksection{V. Meeting Room Casting in Practice — No More Panic}

Meetings are the most frequent casting scenario and also the most prone to problems. Here are practical tips:

\begin{itemize}
  \item \textbf{Arrive Early to Test}: Never wait until the meeting starts to connect casting. Arrive 10 minutes early and test your connection method. Wi-Fi casting fails more often than you'd think.
  \item \textbf{Prepare a Wired Plan B}: Keep an HDMI adapter cable in your bag. When wireless casting won't connect, a wired connection is your lifesaver.
  \item \textbf{Disable Notification Preview}: When casting, your WeChat messages appear on the meeting room big screen. Before casting, turn off "Lock Screen Notification Content Preview" in settings, or enable "Presentation Mode" (some phones support this — automatically hides notifications during casting).
  \item \textbf{"Present Documents" Rather Than "Operate Your Phone"}: If you're presenting a PPT, store the file on your phone and open it in WPS for fullscreen playback, rather than casting your entire phone screen. This is more professional and smoother.
\end{itemize}

\begin{tcbpractice}{Hands-On: Complete a Phone-to-TV Cast}
\begin{enumerate}
  \item Chapter 18: Data Backup and Migration
  \item Open a photo or video on your phone.
  \item Try wireless casting (iPhone uses Screen Mirroring, Android uses the wireless cast button).
  \item If it fails, check whether the TV supports the corresponding protocol, or try DLNA push (tap the "Cast" button in the Gallery/Video app).
  \item After successful casting, try doing other operations on your phone to see if the big screen mirrors them.
\end{enumerate}

\end{tcbpractice}

\begin{tcbsummary}{Chapter Summary}
\begin{itemize}
  \item Three wireless casting camps: AirPlay (Apple), Miracast (Android), DLNA (cross-platform media push)
  \item Wired casting is more stable with zero latency — an HDMI adapter cable is essential for meeting rooms
  \item Samsung DeX / Huawei Desktop Mode turns your phone into a light office computer
  \item Car connectivity: Four major solutions — CarPlay / Android Auto / HiCar / CarLife
  \item Meeting room casting golden rules: test early, prepare a Plan B, disable notification preview
\end{itemize}

\end{tcbsummary}
